训练预算批下来了，一个固定的计算量，用完为止。现在要在两张纸上各填一个数：模型做多大，训多少数据。填完之后模型会训出来，之后还要被四十个客户各自改造一遍，再压进一张放得下它的卡里上线。这一串决定里没有一个能靠直觉拍板，而它们背后其实只有两个数学对象——一条经验幂律，和一个约束优化问题。

这一单元的贯穿线索是：规模带来的收益是可以预测的幂律，但幂律的指数很小，所以每一次等量的能力提升，代价都在按指数上涨；于是``怎样不重新训练就改变模型的行为''成了与``把模型做多大''同等重要的问题。前两节处理规模侧，后两节处理适配与交付侧，它们不是两个话题，是同一个问题的两端。

第一节把幂律的指数换算成人能理解的倍数。\(\alpha \approx 0.07\) 意味着损失减半需要把规模放大约两万倍，``再大十倍''只买到百分之十几的下降。对数-对数纸上那条平缓的直线之所以骗人，是因为横轴每往右一格就是十倍成本。这一节还要把三件事分开：幂律是拟合出来的而非推出来的；\(L_\infty\) 那一项是信息论给的地板，与算力无关；损失的平滑下降与下游能力的关系至今没有定论。

第二节把预算分配写成约束优化。计算量近似是 \(6ND\)，损失是两条幂律之和，一次 Lagrange 乘子就给出 \(N^\star \propto C^{\beta/(\alpha+\beta)}\) 与 \(D^\star \propto C^{\alpha/(\alpha+\beta)}\)——当两个指数接近时，参数与数据应当同比例增长。这个纯粹由数学给出的结论说明早期那批大模型普遍训练数据不足。但最优点属于目标函数而非模型：把推理成本也算进去，最优解会明显偏向``更小的模型训练更久''。

第三节转向已经训好的模型。冻结 \(W_0\)、只学一个秩为 \(r\) 的增量 \(BA\)，参数量从 \(dk\) 降到 \(r(d+k)\)，典型取值下不到百分之一。这里要把它和 SVD 的低秩逼近仔细区分开：Eckart--Young 处理的是对已知矩阵的最优逼近，LoRA 是在低秩流形上做优化，两者的最优性不能互相搬运。支撑它的``微调增量本身低秩''是经验观察。

第四节算交付的账。量化改变数值表示而不改变函数，误差方差正比于步长的平方，真正的难点是激活里的离群值把缩放因子撑坏；蒸馏则换掉函数，让小模型去拟合教师的输出分布，等价于最小化一个 KL 散度，软标签比硬标签多带了类别之间的相似结构。两条路径不在同一条成本--质量曲线上。收束处是一条必须写下来的边界：评测集上没掉，不等于没退化。

\subsection{怎么读这一章}

要做规模决策就从第一节顺读到第二节，这两节合起来是一次完整的预算论证。手上已经有模型、要解决的是改造或部署问题，可以直接从第三节进入。四节都在反复做同一件事——把经验规律和定理分开标注——如果只带走一样东西，带走这个习惯。

\subsection{本章篇目}

\begin{itemize}
\tightlist
\item \hyperref[sec:14-Deep-Learning/09-Scaling-and-Adaptation/01]{``损失随规模下降，但只是幂律''}；
\item \hyperref[sec:14-Deep-Learning/09-Scaling-and-Adaptation/02]{``给定预算，参数与数据怎么分''}；
\item \hyperref[sec:14-Deep-Learning/09-Scaling-and-Adaptation/03]{``LoRA 用低秩把微调成本压下来''}；
\item \hyperref[sec:14-Deep-Learning/09-Scaling-and-Adaptation/04]{``量化与蒸馏把成本换成精度''}。
\end{itemize}

四节读完，你能算清楚一笔训练预算该怎么分、一次微调该花多少、一次压缩会赔上什么，但不会得到任何关于``再大多少就能做到什么''的承诺——那种承诺在这条幂律上不存在。
